% preface.tex
\chapter*{Preface}
\addcontentsline{toc}{chapter}{Preface}

\begin{english}
The intersection of Artificial Intelligence (AI) and Fundamental Physics marks one of the most exciting frontiers in modern science. Traditionally, physics has relied on elegant mathematical models and rigorous experimentation. Today, the explosion of data from experiments like the Large Hadron Collider (LHC) and the Laser Interferometer Gravitational-Wave Observatory (LIGO) requires a new set of tools. This book, \textbf{"AI Python for Fundamental Physics"}, is designed to bridge the gap between classical theoretical physics and the burgeoning field of computational intelligence.

Throughout these pages, we explore how machine learning models can solve complex differential equations, classify quantum states, and discover hidden physical laws from raw data. Our goal is to empower graduate students, researchers, and educators with the practical skills needed to deploy AI in their research workflows using the Python ecosystem.
\end{english}

\vspace{1cm}

\begin{thai}
\textbf{คำนำ}

จุดตัดระหว่างปัญญาประดิษฐ์ (AI) และฟิสิกส์พื้นฐานถือเป็นหนึ่งในพรมแดนที่น่าตื่นเต้นที่สุดในวิทยาศาสตร์สมัยใหม่ ตามเนื้อผ้า ฟิสิกส์อาศัยแบบจำลองทางคณิตศาสตร์ที่สง่างามและการทดลองที่เข้มงวด ในปัจจุบัน การระเบิดของข้อมูลจากการทดลอง เช่น Large Hadron Collider (LHC) และ Laser Interferometer Gravitational-Wave Observatory (LIGO) จำเป็นต้องมีชุดเครื่องมือใหม่ หนังสือเล่มนี้ \textbf{"AI Python for Fundamental Physics"} ได้รับการออกแบบมาเพื่อเชื่อมช่องว่างระหว่างฟิสิกส์ทฤษฎีคลาสสิกและสาขาความฉลาดทางคอมพิวเตอร์ที่กำลังเติบโต

ในแต่ละบท เราจะสำรวจว่าโมเดลการเรียนรู้ของเครื่องสามารถแก้สมการเชิงอนุพันธ์ที่ซับซ้อน จำแนกสถานะควอนตัม และค้นพบกฎทางฟิสิกส์ที่ซ่อนอยู่จากข้อมูลดิบได้อย่างไร เป้าหมายของเราคือการเพิ่มขีดความสามารถให้กับนักศึกษาระดับบัณฑิตศึกษา นักวิจัย และอาจารย์ ด้วยทักษะพื้นฐานที่จำเป็นในการนำ AI ไปใช้ในกระบวนการวิจัยโดยใช้ระบบนิเวศของ Python
\end{thai}

\vfill
\noindent \textbf{Asst. Prof. Dr. Chewa Thassana} \\
Rambhai Barni Rajabhat University \\
\today
